\begin{afterword}
\summaryhead

The post asks whether a giant lookup table, a machine that stores a human's response to every possible input, would be a zombie: something that behaves like a human, even writing philosophy papers about consciousness, without being conscious. No one, the author says, claims such a table is conscious.

The author's method is ``Follow-The-Improbability'': a table that matches a human is so improbable that one must ask where it came from. If it was computed from a specification of a brain, it relays a conscious source, as a cellphone does. If it was drawn at random and then picked out by design, the picker is the source. If the philosopher stipulates that it was drawn by pure chance, the improbability lies in the philosopher's specification. The moral is that tracing talk about consciousness back generally finds consciousness. A table drawn by genuinely pure chance, the author grants, would not be conscious. A coda applies the same reasoning to claims that an arbitrary AI would be moral.

\responsehead

The thought experiment and its first answer are older than the post. Ned Block introduced the lookup machine in 1981 to argue against behaviorism about intelligence, and Block already made the post's comparison: the machine is ``a conduit for intelligence,'' like a two-way radio, and ``All the intelligence it exhibits is that of its programmers.'' Shannon and McCarthy had raised the idea in 1956. The post names Block only inside Dennett's quotation. What it adds is the method: measure how improbable the table is, and follow that improbability back to a brain specification, a designer or the philosopher's stipulation.

The method answers a narrower question than the one the post asks. The post asks whether the table behaves ``exactly like a human without being conscious,'' and later says that ``The consciousness isn't in the lookup table.'' In that sense the table is a zombie. The cellphone answer is about where the table's talk comes from. That is the sense the Generalized Anti-Zombie Principle needs, and the moral is stated in it, but the change of sense is not marked.

For the last case the post gives its own opinion. For a table drawn by genuinely pure chance, the post says it ``wouldn't be conscious. IMHO,'' and gives as its reason the reductio stated earlier. So the post allows that, if such a table were drawn, something would talk as we do without being conscious; what the post disputes is that such talk ever arises without a conscious source. That is a claim about probability rather than possibility, and the post says openly that it is leaving ``standard philosophical procedure.''

Two steps state as settled what the argument does not need. Quantum randomness is described through many-worlds (``the deterministic result''), but a random table is just as improbable on any interpretation. And the table computed from a brain writes ``because of a conscious algorithm,'' which assumes what the previous post left for later, that a brain's specification run elsewhere is conscious. The talk would trace back to the human whose brain was specified in any case.

The coda applies the same reasoning to AI design, and says openly why the author cares about it.

In short: Block's lookup table, with Block's own answer that it is a conduit, extended by a method of tracing improbability to its source; it shows that the talk of any table we could meet would come from consciousness somewhere, and grants, as opinion, that a table made by pure chance would talk without being conscious.
\end{afterword}
